% =====================================================================
\chapter{Logistic Regression}
% =====================================================================

\labinfo{06\_Logistic\_Regression}{logistic\_regression.py}{Breast Cancer Wisconsin (Diagnostic): 569
tumours, 30 features}{Binary classification (malignant or benign)}{Test accuracy 97.3\%, recall 97.6\%,
ROC AUC 0.997}

\section{Dataset}
The \textbf{Breast Cancer Wisconsin (Diagnostic)} dataset (UCI) describes 569 breast tumours, 212
malignant (37.3\%) and 357 benign. The features come from a digitised microscope image of cells taken by
fine-needle aspiration. Ten properties of the cell nuclei are measured (radius, texture, perimeter, area,
smoothness, compactness, concavity, concave points, symmetry and fractal dimension), and for each one the
dataset gives the mean, the standard error and the ``worst'' value (mean of the three largest), so 30
features in total. They are on very different scales (area is in the hundreds, smoothness around
0.1), so I standardised them. I used a stratified 80/20 split: 456 training and 113 test tumours.
Malignant is the positive class.

\section{Code and experiments}
\begin{itemize}
\item \code{sigmoid} clips its input to $[-500,500]$ so \code{np.exp} never overflows, and
      \code{cross\_entropy} clips probabilities away from 0 and 1 before taking logs.
\item \code{LogisticRegression.fit} runs batch gradient descent (learning rate 0.1, 2,000 iterations)
      and stores the loss after each iteration; \code{predict\_proba} and \code{predict} use the
      learned weights.
\item \code{binary\_metrics} and \code{roc\_auc} compute the metrics from scratch.
\end{itemize}
\noindent Experiments: (1)~5-fold cross-validation for four L2 strengths; (2)~the chosen model on the
test set at thresholds 0.5 and 0.3; (3)~a two-feature version of the model, only to draw the probability
surface.

\section{Results}

\begin{table}[H]
\centering\small
\caption{5-fold cross-validation on the training set.}
\label{tab:logreg-cv}
\begin{tabular}{@{}l c c c c@{}}
\toprule
L2 strength $\lambda$ & 0 & 0.001 & 0.01 & 0.1\\
\midrule
CV accuracy & $0.987\pm0.011$ & $0.987\pm0.011$ & $0.985\pm0.009$ & $0.969\pm0.008$\\
CV ROC AUC & 0.9929 & 0.9930 & 0.9928 & 0.9921\\
\bottomrule
\end{tabular}
\end{table}

Regularisation did not help: with 456 examples for 31 weights and early stopping after 2,000 steps,
the plain maximum-likelihood model was already as good as the regularised ones. I kept $\lambda=0$.

\begin{table}[H]
\centering\small
\caption{Test-set results (113 tumours, 42 malignant).}
\label{tab:logreg-test}
\begin{tabular}{@{}l S[table-format=1.4] S[table-format=1.4] S[table-format=1.4] S[table-format=1.4] S[table-format=1.4] l@{}}
\toprule
\textbf{Threshold} & {\textbf{Accuracy}} & {\textbf{Precision}} & {\textbf{Recall}} & {\textbf{Specificity}} & {\textbf{F1}} & \textbf{[[TN FP] [FN TP]]}\\
\midrule
0.5 & 0.9735 & 0.9535 & 0.9762 & 0.9718 & 0.9647 & [[69 2] [1 41]]\\
0.3 & 0.9558 & 0.9111 & 0.9762 & 0.9437 & 0.9425 & [[67 4] [1 41]]\\
\midrule
\multicolumn{7}{@{}l}{ROC AUC = 0.9973 \qquad test log-loss (cross-entropy) = 0.0752}\\
\bottomrule
\end{tabular}
\end{table}

The model gets 110 of 113 tumours right and ranks almost every malignant tumour above every benign one
(AUC 0.997). Lowering the threshold to 0.3 did \emph{not} help on this test set: the one missed malignant
tumour received a probability of only 0.06, so it stays missed, and the lower threshold just adds two more
false alarms. To catch it, the threshold would have to go below 0.06, which would flag many healthy
patients.

The largest weights (on standardised features) belong to worst texture ($+1.17$), standard error of the
radius ($+1.17$), worst smoothness ($+1.03$), worst radius ($+0.99$) and mean concave points ($+0.98$).
All are positive: large, irregular nuclei push the prediction towards malignant, which agrees with
what is known about cancer cells.

\begin{figure}[H]
\centering
\includegraphics[width=0.48\linewidth]{\figpath{06\_Logistic\_Regression}{loss_and_roc.pdf}}\hfill
\includegraphics[width=0.48\linewidth]{\figpath{06\_Logistic\_Regression}{threshold_effect.pdf}}
\caption{Training loss, ROC curve and threshold trade-off.}
\end{figure}

\section{Source code}
\lstinputlisting[caption={\texttt{06\_Logistic\_Regression/logistic\_regression.py}}]{06_Logistic_Regression/logistic_regression.py}
\clearpage
